\SourcePage{245}
\SGAsection{4}{一般情形}
我们将概述证明 2.5 的另一种方法，它适用于任意\Term{characteristic}{特征}；这里承认下面的断言。它与 3.1 相近，是第 XV 讲 1.1.4 的特例。

\begin{SGAresult}{断言}{4.1}
$\mathbf P(N)$ 中由满足下述条件的点 $(x,H)$ 组成的\Term{subset}{子集}是\Term{open-subset}{开集}：$x$ 是 $X\cdot H$ 的一个\Term{ordinary-quadratic-singularity}{普通二次奇点}。
\end{SGAresult}
注意，除 $\operatorname{car}(k)=2$ 且 $\dim X$ 为奇数的情形外，4.1 都由 3.3 \Term{pushout}{推出}。

\textbf{4.1.1。}以 $R_1$ 表示 $\mathbf P(N)$ 中与 4.1 的\Term{open-subset}{开集}互补的\Term{closed-subset}{闭集}，以 $F_1$ 表示它的像
\[
\varphi(R_1)\subset\check X\subset\check{\mathbf P}^{r}.
\]
\begin{SGAresult}{\Term{lemma}{引理}}{4.1.2}
若 4.1 的\Term{open-subset}{开集}\Term{nonempty}{非空}，则
$\operatorname{codim}_{\check{\mathbf P}^{r}}(F_1)\geq2$。
\end{SGAresult}
\begin{proof}
由假设，$\dim(R_1)<\dim(\mathbf P(N))=r-1$，而 $F_1=\varphi(R_1)$。
\end{proof}

\textbf{4.1.3。}以 $F_2$ 表示 $\check X$ 中由在至少两个点与 $X$ \Term{tangent}{相切}的\Term{hyperplane}{超平面}组成的部分。注意，$F_2$ 并非\Term{closed-subset}{闭集}，但有
\SourcePage{246}
\[
\overline{F_2}\subset F_1\cup F_2.
\tag{4.1.4}
\]
这是因为 $F_1\cup F_2=F$（参见 3.2）是\Term{closed-subset}{闭集}。（包含关系 4.1.4 “表明”，两个\Term{singularity}{奇点}汇合而成的\Term{singularity}{奇点}不是\Term{ordinary-quadratic-singularity}{普通二次奇点}。）无论如何，4.2.4 的构造将表明 $F_2$ 是\Term{constructible}{可构造的}。为证明 2.5.1，我们将应用判据 3.2.4。为此，需要知道
\[
\operatorname{codim}_{\check{\mathbf P}^{r}}(F_1)\geq2,
\qquad
\operatorname{codim}_{\check{\mathbf P}^{r}}(F_2)\geq2.
\]
另一方面，3.8.1 和 4.1.2 保证，将给定的\Term{embedding}{嵌入}换成它的 $d$ 倍（$d\geq2$）之后，将有
$\operatorname{codim}_{\check{\mathbf P}^{r}}(F_2)\geq2$。因此，只要知道下面的命题，就证明了 2.5.1。

\begin{SGAresult}{命题}{4.2}
设 $X$ 是 $\mathbf P^r$ 的一个\Term{smooth}{光滑}\Term{irreducible}{不可约}\Term{subvariety}{子簇}，\Term{dimension}{维数}为 $n\geq1$。对任意整数 $d\geq2$，在
\[
\check{\mathbf P}^{\binom{r+d}{d}-1}=\check{\mathbf P}^{N}
\]
中，由 $\mathbf P^r$ 中在至少两个点与 $X$ \Term{tangent}{相切}的 $d$ 次\Term{hypersurface}{超曲面}组成的部分 $F_2$，在 $\check{\mathbf P}^{N}$ 中的\Term{codimension}{余维}至少为 $2$。
\end{SGAresult}

\textbf{4.2.1。}先考虑 $d=2$ 且 $X$ 是 $\mathbf P^r$ 的一个\Term{linear-subvariety}{线性子簇}的特殊情形。只须在 $\check{\mathbf P}^{N}$ 中找一条不与 $F$ 相交的直线即可（参见 3.2.4）。在 $\mathbf P^r$ 中选取\Term{homogeneous-coordinate}{齐次坐标} $X_0,\ldots,X_r$，使得 $X$ 由
\SourcePage{247}
$X_{n+1}=\cdots=X_r=0$ 定义，并在 $\mathbf P^1$ 中选取\Term{homogeneous-coordinate}{齐次坐标} $(\lambda,\mu)$。

当 $n$ 为偶数，$n=2k$ 时，取如下\Term{quadric}{二次曲面}的\Term{pencil}{铅笔}：
\[
H_{(\lambda,\mu)}=
\lambda X_0^2+(\lambda X_k+\mu X_0)X_{2k}
+\sum_{i=1}^{k-1}(\lambda X_i+\mu X_{i+1})X_{k+i};
\tag{4.2.2}
\]
当 $n$ 为奇数，$n=2k+1$ 时，取\Term{pencil}{铅笔}
\[
H_{(\lambda,\mu)}=(\lambda X_1+\mu X_0)X_0
+\sum_{i=1}^{k}(\lambda X_{2i}+\mu X_i)X_{2i+1}.
\tag{4.2.3}
\]
在这两种情形中，都可看出：当 $(\lambda,\mu)\ne(0,1)$ 时，$X\cdot H_{(\lambda,\mu)}$ 是\Term{smooth}{光滑的}；而 $X\cdot H_{(0,1)}$ 只有一个\Term{singularity}{奇点}，且它还是一个\Term{ordinary-quadratic-singularity}{普通二次奇点}。所以该\Term{pencil}{铅笔}不与 $F$ 相交。

\textbf{4.2.4。}以 $\operatorname{Btg}^{d}(X)$ 表示
$\check{\mathbf P}^{N}\times(X\times X-\Delta(X\times X))$ 中由三元组 $(H,x_1,x_2)$ 组成的\Term{subvariety}{子簇}，其中 $H$ 是 $\mathbf P^r$ 中的一个 $d$ 次\Term{hypersurface}{超曲面}，在 $x_1$ 和 $x_2$ 与 $X$ \Term{tangent}{相切}，且 $x_1\ne x_2$。$\operatorname{Btg}^{d}(X)$ 在 $\operatorname{pr}_1$ 下的像就是（4.2）中的\Term{set}{集合} $F_2$，因此只须证明
\[
\dim\operatorname{Btg}^{d}(X)\leq N-2
\quad\text{若 }d\geq3,\ \text{或 }d=2\text{ 且 }X\text{ 不是线性的}。
\tag{4.2.5}
\]
为此，以 $\operatorname{Lin}(X\times X)$ 表示 $X\times X-\Delta(X\times X)$ 中的\Term{subvariety}{子簇}
\[
\begin{aligned}
\operatorname{Lin}(X\times X)
&=\{(x_1,x_2)\in X\times X\mid x_1\ne x_2,\ \text{且由 }x_1,x_2
\text{ 张成的直线在 }x_1\text{ 与 }X\text{ 相切}\}\\
&=\{(x_1,x_2)\in X\times X\mid x_1\ne x_2,\ x_2\in X\cap\operatorname{Tg}_X(x_1)\}.
\end{aligned}
\tag{4.2.6}
\]
\SourcePage{248}
显然有：
\begin{SGAresult}{断言}{4.2.7}
若 $X$ 不是 $\mathbf P^r$ 的\Term{linear-subvariety}{线性子簇}，则
\[
\dim\operatorname{Lin}(X\times X)<2\dim X.
\]
\end{SGAresult}
因此，为证明（4.2.5），只须证明\Term{morphism}{态射}
\[
\operatorname{Btg}^{d}(X)\xrightarrow{\operatorname{pr}_{2,3}}
X\times X-\Delta(X\times X)
\tag{4.2.8}
\]
的\Term{fiber}{纤维}，视为 $\check{\mathbf P}^{N}$ 的\Term{subvariety}{子簇}时，满足
\[
\begin{cases}
\text{在 }\operatorname{Lin}(X\times X)\text{ 上方的余维}\geq2\dim(X)+1,\\
\text{在 }X\times X-\Delta(X\times X)-\operatorname{Lin}(X\times X)
\text{ 上方的余维}\geq2\dim(X)+2.
\end{cases}
\tag{4.2.9}
\]
在下面的\Term{lemma}{引理}中取 $L_i$ 为 $X$ 在 $x_i$ 的\Term{tangent-space}{切空间}（$i=1,2$），便得到断言（4.2.9）。

\begin{SGAresult}{命题}{4.3}
设 $x_1,x_2$ 是 $\mathbf P^r$ 的两个相异\Term{closed-point}{闭点}，$L_1,L_2$ 是 $\mathbf P^r$ 的两个\Term{linear-subspace}{线性子空间}，且 $x_1\in L_1$、$x_2\in L_2$。在 $d$ 次\Term{hypersurface}{超曲面}的\Term{space}{空间}
$\check{\mathbf P}^{N}=\check{\mathbf P}^{\binom{r+d}{d}-1}$ 中，由在 $x_1$ 与 $L_1$ \Term{tangent}{相切}并在 $x_2$ 与 $L_2$ \Term{tangent}{相切}的\Term{hypersurface}{超曲面}组成的\Term{linear-subspace}{线性子空间}，其\Term{codimension}{余维}等于
\[
\begin{cases}
\dim L_1+\dim L_2+2,&
\begin{array}{l}d\geq3,\text{ 或 }d=2\text{ 且 }L_1\cap L_2\\
\text{不包含经过 }x_1,x_2\text{ 的直线 }(x_1,x_2),\end{array}\\[2pt]
\dim L_1+\dim L_2+1,&d=2\text{ 且 }L_1\cap L_2\text{ 包含直线 }(x_1,x_2).
\end{cases}
\]
\end{SGAresult}
\SourcePage{249}
\begin{proof}
当 $d\geq3$ 时，要证明 $\dim L_1+\dim L_2+2$ 个\Term{linear}{线性}条件相互独立，因而只须处理 $L_1=L_2=\mathbf P^r$ 的情形，即求出“有多少”\Term{hypersurface}{超曲面}在 $x_1$ 和 $x_2$ 都是\Term{singular}{奇异的}。选取\Term{homogeneous-coordinate}{齐次坐标}，使得
$x_1=(1,0,\ldots,0)$、$x_2=(0,1,0,\ldots,0)$。方程为 $\sum A_wX^w$ 的\Term{hypersurface}{超曲面}在 $(1,0,\ldots,0)$ 是\Term{singular}{奇异的}，当且仅当\Term{monomial}{单项式}
$X_0^{d-1}X_i$（$i=0,\ldots,r$）的系数全为零；它在 $(0,1,0,\ldots,0)$ 是\Term{singular}{奇异的}，当且仅当\Term{monomial}{单项式} $X_1^{d-1}X_i$（$i=0,\ldots,r$）的系数全为零。当 $d\geq3$ 时，\Term{monomial}{单项式}
$X_0^{d-1}X_i$（$i=0,\ldots,r$）与 $X_1^{d-1}X_i$（$i=0,\ldots,r$）彼此都不同，这就给出了 $d\geq3$ 时所需的断言。

当 $d=2$ 时，同一计算给出不等式
\[
\operatorname{codim}\geq\dim L_1+\dim L_2+1
\]
（因为\Term{monomial}{单项式} $X_0X_1$ 被计数了两次）。

现在验证，若 $L_1\cap L_2$ 不包含直线 $(x_1,x_2)$，则\Term{tangent}{相切}条件相互独立。可归结为 $L_1=\mathbf P^r$，而 $L_2$ 是 $\mathbf P^r$ 中不包含直线 $(x_1,x_2)$ 的一个\Term{hyperplane}{超平面}的情形。选取\Term{homogeneous-coordinate}{齐次坐标} $X_0,\ldots,X_r$，使得
$x_1=(1,0,\ldots,0)$、$x_1=(0,1,0,\ldots,0)$，且 $L_2$ 是\Term{hyperplane}{超平面} $X_0=0$。\Term{quadric}{二次曲面} $\sum_{i\leq j}A_{ij}X_iX_j$ 在 $(1,0,\ldots,0)$ 是\Term{singular}{奇异的}，当且仅当 $A_{0j}=0$（$j=0,\ldots,r$）；它在 $(0,1,\ldots,0)$ 与
\SourcePage{250}
$X_0=0$ \Term{tangent}{相切}，当且仅当 $A_{1j}=0$（$j=1,\ldots,r$）。合起来得到 $2r+1=\dim L_1+\dim L_2+2$ 个相互独立的条件。

对于 $L_1\cap L_2$ 包含直线 $(x_1,x_2)$ 的情形，可归结为 $L_1=L_2=\text{直线 }(x_1,x_2)$ 的情形。在使得 $x_1=(1,0,\ldots,0)$、$x_2=(0,1,0,\ldots,0)$ 的\Term{homogeneous-coordinate}{齐次坐标}中，\Term{quadric}{二次曲面} $\sum_{i\leq j}A_{ij}X_iX_j$ 经过 $(1,0,\ldots,0)$ 且在该点与 $L_1$ \Term{tangent}{相切}，当且仅当 $A_{00}=A_{01}=0$；它经过 $(0,1,0,\ldots,0)$ 且在该点与 $L_2$ \Term{tangent}{相切}，当且仅当 $A_{11}=A_{0,1}=0$。合起来只有 $3=\dim L_1+\dim L_2+1$ 个条件。
\end{proof}

\SGAsection{5}{用“初等”方法求对偶簇的次数}
（当 $\dim X$ 为偶数或 $\operatorname{car}(k)\ne2$ 时，另一种方法见第 XVIII 讲 3.2。）

我们仍沿用 2.2 的记号，因此 $X\hookrightarrow\mathbf P^r$ 是一个\Term{dimension}{维数}为 $n$ 的\Term{smooth}{光滑}\Term{irreducible}{不可约}\Term{subvariety}{子簇}。

\textbf{5.1。}先从更具“\Term{functorial}{函子性}”的角度来考察\Term{morphism}{态射}（3.1.1）
\[
\nu:\mathbf P(N)\longrightarrow X\times\check{\mathbf P}^{r}.
\tag{5.1.0}
\]
（我要感谢 S. Kleiman，他向我指出了下面的表述。）在 $X$ 上有一个\Term{coherent-sheaf}{凝聚层}的\Term{exact-sequence}{正合列}
\SourcePage{251}
\[
0\longrightarrow N^\vee\longrightarrow
\Omega^1_{\mathbf P/k}|_X\longrightarrow\Omega^1_X\longrightarrow0.
\tag{5.1.1}
\]
另一方面，“用\Term{homogeneous-coordinate}{齐次坐标}写出 $\mathbf P^r$ 上的\Term{differential-form}{微分形式}”，便得到 $\mathbf P^r$ 上的\Term{sheaf}{层}的\Term{exact-sequence}{正合列}
\[
0\longrightarrow\Omega^1_{\mathbf P^r/k}
\longrightarrow\bigl(\mathcal O_{\mathbf P}(-1)\bigr)^{r+1}
\longrightarrow\mathcal O_{\mathbf P}\longrightarrow0;
\tag{5.1.2}
\]
将它\Term{restriction}{限制}到 $X$ 上，便得到\Term{exact-sequence}{正合列}
\[
0\longrightarrow\Omega^1_{\mathbf P^r/k}|_X
\longrightarrow\bigl(\mathcal O_X(-1)\bigr)^{r+1}
\longrightarrow\mathcal O_{\mathbf P}\longrightarrow0.
\tag{5.1.3}
\]
结合 5.1.1 和 5.1.3，得到一个\Term{inclusion}{包含映射}
\[
N^\vee\lhook\joinrel\longrightarrow
\bigl(\mathcal O_X(-1)\bigr)^{r+1},
\tag{5.1.4}
\]
取\Term{transpose}{转置}便给出一个\Term{surjection}{满射}
\[
\mathcal O_X(1)^{r+1}\longrightarrow N,
\tag{5.1.5}
\]
从而得到一个\Term{closed-immersion}{闭浸入}
\[
\alpha:\mathbf P(N)\longrightarrow
\mathbf P\bigl(\mathcal O_X(1)^{r+1}\bigr)
\simeq X\times\check{\mathbf P}^{r},
\tag{5.1.6}
\]
它就是（5.1.0）。

\SourcePage{252}
\textbf{5.2。计算的思路。}

从现在起，假设 $\check X=\varphi(\mathbf P(N))$ 的\Term{dimension}{维数}为 $r-1$，因而它是 $\check{\mathbf P}^{r}$ 中的一个\Term{hypersurface}{超曲面}。引入\Term{chern-class}{Chern 类} [1]，就可以自然地计算它的\Term{degree}{次数}。

对任意\Term{projective}{射影}\Term{smooth}{光滑}\Term{variety}{簇} $T$，以 $\operatorname{CH}(T)$ 表示它的\Term{chow-ring}{Chow 环}；这是一个\Term{graded-ring}{分次环}（按\Term{cycle}{圈}的\Term{codimension}{余维}\Term{graded}{分次}）。给定元素 $z\in\operatorname{CH}(T)$，以
\[
\int_T z
\]
表示它的零次\Term{set-component}{分量}（一个 $0$-\Term{cycle}{圈}）的\Term{degree}{次数}。当不致混淆时（例如在积分号 $\int_T$ 下），$T$ 上\Term{invertible-sheaf}{可逆层} $L$ 在 $\operatorname{CH}^1(T)$ 中的类也记为 $L$。

于是，以 $L$ 表示一条直线在 $\operatorname{CH}(\check{\mathbf P}^{r})$ 中的类，有
\[
\deg(\check X)=\int_{\check{\mathbf P}^{r}}L\cdot\check X.
\tag{5.2.1}
\]
而从“\Term{set}{集合}意义”看，$\check X$ 是 $\mathbf P(N)$ 在 $\varphi$ 下的像（3.1.2）。但对于“\Term{direct-image}{直像}”\Term{homomorphism}{同态}
\[
\varphi_*:\operatorname{CH}(\mathbf P(N))
\longrightarrow\operatorname{CH}(\check{\mathbf P}^{r}),
\]
\SourcePage{253}
有
\[
\varphi_*(\mathbf P(N))=\deg(\varphi)\cdot\check X,
\]
所以 5.2.1 可改写为
\[
\deg(\check X)=\frac1{\deg(\varphi)}
\int_{\check{\mathbf P}^{r}}L\cdot\varphi_*(\mathbf P(N)).
\tag{5.2.2}
\]
这里略微滥用记号：在 $\deg(\varphi)$ 中，$\varphi$ 显然指由（3.1.2）诱导的\Term{morphism}{态射} $\mathbf P(N)\longrightarrow\check X$；其\Term{degree}{次数}是一个非负整数，且它为零当且仅当 $\dim\check X<r-1=\dim(\mathbf P(N))$。

由 $\varphi$ 的\Term{projection-formula}{投影公式}，上式又可改写为
\[
\deg(\check X)=\frac1{\deg(\varphi)}
\int_{\mathbf P(N)}\varphi^*(L).
\tag{5.2.3}
\]
以 $\eta\in\operatorname{CH}(\check{\mathbf P}^{r})$ 表示一个\Term{hyperplane}{超平面}的类，于是 $L=\eta^{r-1}$，而（5.2.3）改写为
\[
\deg(\check X)=\frac1{\deg(\varphi)}
\int_{\mathbf P(N)}\bigl(\varphi^*(\eta)\bigr)^{r-1}.
\tag{5.2.4}
\]
现在用“具体”的方式计算类 $\varphi^*(\eta)$。以 $\xi$（分别以 $\tau$）表示 $\mathbf P(N)$ 上的\Term{canonical}{典范}\Term{invertible-sheaf}{可逆层} $\mathcal O_{\mathbf P(N)}(1)$（分别表示 $\mathbf P(\mathcal O_X(1)^{r+1})$ 上的 $\mathcal O_{\mathbf P(\mathcal O_X(1)^{r+1})}$）。由\Term{functorial}{函子性}，有（参见（5.1.6））
\[
\alpha^*(\tau)=\xi.
\tag{5.2.5}
\]
\SourcePage{254}
另一方面，借助等同
$\mathbf P(\mathcal O_X(1)^{r+1})\simeq X\times\check{\mathbf P}^{r}$，置 $H=\mathcal O_X(1)$，显然有
\[
\tau=\operatorname{pr}_1^*(H)\otimes\operatorname{pr}_2^*(\eta).
\tag{5.2.6}
\]
以 $\pi:\mathbf P(N)\longrightarrow X$ 表示\Term{structure-morphism}{结构态射}，有
\[
\begin{cases}
\pi=\operatorname{pr}_1\circ\alpha,\\
\varphi=\operatorname{pr}_2\circ\alpha.
\end{cases}
\]
因此，对（5.2.6）的两边\Term{action}{作用} $\alpha^*$，由（5.2.5）得到
\[
\xi=\alpha^*(\tau)=\pi^*(H)\otimes\varphi^*(\eta).
\tag{5.2.6}
\]
所以在 $\operatorname{CH}(\mathbf P(N))$ 中有
\[
\varphi^*(\eta)=\xi-\pi^*(H),
\tag{5.2.7}
\]
而（5.2.3）改写为
\[
\deg(\check X)=\frac1{\deg(\varphi)}
\int_{\mathbf P(N)}\bigl(\xi-\pi^*(H)\bigr)^{r-1}.
\tag{5.2.8}
\]
借助下面的\Term{lemma}{引理}（参见 5.3.3），便能“计算”5.2.8 的右端。

\begin{SGAresult}{\Term{lemma}{引理}}{5.3}
设 $X$ 是\Term{dimension}{维数}为 $n$ 的\Term{quasi-projective}{拟射影}\Term{smooth}{光滑}\Term{connected}{连通}\Term{variety}{簇}，$M$ 是 $X$ 上一个\Term{finite-rank}{有限秩} $a$ 的\Term{locally-free-sheaf}{局部自由层}。以 $\xi$ 表示 $\mathbf P(M)$ 上的\Term{invertible-sheaf}{可逆层} $\mathcal O_{\mathbf P(M)}(1)$，
\SourcePage{255}
以 $\pi:\mathbf P(M)\longrightarrow X$ 表示\Term{structure-morphism}{结构态射}。已知 $\operatorname{CH}(\mathbf P(M))$ 是一个（经由 $\pi^*$ 的）\Term{free-module}{自由} $\operatorname{CH}(X)$-\Term{module}{模}，以 $1,\xi,\ldots,\xi^{a-1}$ 为基，且 $\xi$ 满足方程
\[
\xi^a-c_1\xi^{a-1}+c_2\xi^{a-2}-c_3\xi^{a-3}\cdots=0,
\tag{5.3.1}
\]
其中 $c_i$ 表示 $M$ 在 $\operatorname{CH}(X)$ 中的\Term{chern-class}{Chern 类}。置
$C(M)=1+c_1+c_2+\cdots$，即 $M$ 的\Term{total-chern-class}{总 Chern 类}；它是 $\operatorname{CH}(X)$ 中的\Term{invertible-element}{可逆元}。于是：

\textup{(i)} 有公式
\[
\pi_*\left(\frac{\xi^{a-1}}{1+\xi}\right)=\frac1{C(M)}.
\tag{5.3.2}
\]

\textup{(ii)} 设 $F(T)\in\operatorname{CH}(X)[T]$ 是系数在 $\operatorname{CH}(X)$ 中的\Term{polynomial}{多项式}，且 $T^i$ 的系数是\Term{degree}{次数}为 $n+a-1-i$ 的\Term{homogeneous-element}{齐次元}。因而 $F(\xi)$ 是 $\operatorname{CH}(\mathbf P(M))$ 中\Term{degree}{次数}为“最大值”$n+a-1=\dim\mathbf P(M)$ 的元素。在这些条件下，有
\[
\int_{\mathbf P(M)}F(\xi)=(-1)^{a-1}\int_X\frac{F(-1)}{C(M)}.
\tag{5.3.3}
\]
\end{SGAresult}
\begin{proof}
先由（5.3.2）\Term{pushout}{推出}（5.3.3）。写成
\[
\frac1{C(M)}=\sum_{j\geq0}b_j,
\qquad
\begin{gathered}
b_0=1,\quad b_j\text{ 是 }\operatorname{CH}(X)\text{ 中次数为 }j\text{ 的元素},\\
b_j=0\quad(j<0),
\end{gathered}
\tag{5.3.4}
\]
\SourcePage{256}
则（5.3.2）改写为
\[
\pi_*(\xi^{a-1+j})=(-1)^jb_j
\qquad(a-1+j\geq0)
\tag{5.3.5}
\]
（注意，当 $j<a-1$ 时，由\Term{degree}{次数}的原因有 $\pi_*(\xi^j)=0$）。由此\Term{pushout}{推出}（5.3.3）。根据\Term{linear}{线性}性，只须处理 $F(T)=uT^{a-1+j}$ 的情形，其中 $u$ 是 $\operatorname{CH}(X)$ 中\Term{degree}{次数}为 $n-j$ 的\Term{homogeneous-element}{齐次元}。由（5.3.5）和 $\pi$ 的\Term{projection-formula}{投影公式}，有
\[
\begin{aligned}
\int_{\mathbf P(M)}F(\xi)
&=\int_{\mathbf P(M)}\pi^*(u)\cdot\xi^{a-1+j}\\
&=\int_Xu\cdot\pi_*(\xi^{a-1+j})
=(-1)^j\int_Xu\cdot b_j.
\end{aligned}
\tag{5.3.6}
\]
另一方面，
\[
\begin{aligned}
(-1)^{a-1}\int_X\frac{F(-1)}{C(M)}
&=(-1)^{a-1}\int_X\frac{(-1)^{a-1+j}u}{C(M)}\\
&=(-1)^j\int_X\sum_i u\,b_i,
\end{aligned}
\tag{5.3.7}
\]
而 $\deg(u)=n-j$、$\deg(b_i)=i$，故
\[
\int_Xu\cdot b_i=0\qquad(i\ne j),
\]
所以（5.3.7）改写为
\[
(-1)^{a-1}\int_X\frac{F(-1)}{C(M)}
=(-1)^j\int_Xu\cdot b_j.
\tag{5.3.8}
\]
这就得到了所说的蕴含关系。

现在证明（5.3.2），或者\Term{equivalence}{等价}地证明（5.3.5）。
\SourcePage{257}
已知由\Term{degree}{次数}的原因，
\[
\pi_*(\xi^j)=0\qquad(j<a-1),
\]
而由几何上的原因（$\xi^{a-1}$ 与每个\Term{fiber}{纤维}“交于一点”），有
\[
\pi_*(\xi^{a-1})=1_X.
\]
因而由\Term{projection-formula}{投影公式}，有
\[
\pi_*\left(\sum_{j=0}^{a-1}\pi^*(U_i)\,\xi^i\right)=U_{j-i}.
\tag{5.3.9}
\]
由（5.3.9），考虑到方程 5.3.1，断言 5.3.5 是下面\Term{lemma}{引理}的推论；在该\Term{lemma}{引理}中，$\xi$ 换成 $T$，$c_i$ 换成 $S_i$。
\end{proof}

\begin{SGAresult}{\Term{lemma}{引理}}{5.4}
设 $a$ 为正整数，以 $A$ 表示含 $a$ 个\Term{indeterminate}{不定元}的\Term{formal-power-series-ring}{形式幂级数环} $\mathbf Z[[X_1,\ldots,X_a]]$。以 $B$ 表示 $A$-\Term{algebra}{代数}
\[
B=A[[T]]\Big/\left(\prod_{i=1}^{a}(T-X_i)\right).
\]
于是 $B$ 是一个\Term{free-module}{自由} $A$-\Term{module}{模}，以 $1,T,\ldots,T^{a-1}$ 为基，
\[
B=AT^{a-1}\oplus AT^{a-2}\oplus\cdots\oplus A.
\tag{5.4.1}
\]
\SourcePage{258}
以 $S_i\in A$ 表示第 $i$ 个\Term{symmetric-function}{对称函数}
\[
S_i=\sum_{1\leq\alpha_1<\cdots<\alpha_i\leq a}
X_{\alpha_1}\cdots X_{\alpha_i},
\]
于是
\[
\sum_{i=0}^{a}(-1)^iS_iT^{a-i}=0\quad\text{在 }B\text{ 中}。
\tag{5.4.2}
\]
元素 $1+S_1+S_2+\cdots+S_a$ 在 $A$ 中是\Term{invertible-element}{可逆元}，将它的逆写成
\[
\frac1{\sum_{i=0}^{a}S_i}=\sum_{j\geq0}b_j,
\qquad b_j\text{ 是 }A\text{ 中次数为 }j\text{ 的齐次元},\quad b_0=1.
\tag{5.4.3}
\]
另一方面，对每个整数 $j\geq0$，根据 5.4.1 写成
\[
T^{a-1+j}=a_jT^{a-1}+\mathord?\,T^{a-2}
+\mathord?\,T^{a-3}+\cdots+\mathord?,
\qquad\text{其中各 }?\in A,
\tag{5.4.4}
\]
由此定义 $A$ 中一个\Term{degree}{次数}为 $j$ 的\Term{homogeneous-element}{齐次元} $a_j$。在这些条件下，对每个 $j\geq0$ 都有
\[
a_j=(-1)^jb_j.
\tag{5.4.5}
\]
\end{SGAresult}
\begin{proof}
考虑到 $b_j$ 的定义（5.4.3），只须证明
\SourcePage{259}
\[
\begin{gathered}
a_0=1,\\
\sum_{i=0}^{a}(-1)^iS_i a_{j-i}=0\qquad(j>0).
\end{gathered}
\tag{5.4.6}
\]
显然 $a_0=1$，而第二个关系（对给定的 $j>0$）正是利用 5.4.1，在关系
\[
\sum_{i=0}^{a}(-1)^iS_iT^{a-1+j-i}=0\quad\text{在 }B\text{ 中}
\tag{5.4.7}
\]
中取 $T^{a-1}$ 的系数所得到的；这个关系由 5.4.2 乘以 $T^{j-1}$ 得到。证毕。
\end{proof}

\begin{SGAresult}{推论}{5.5}
有公式
\[
\deg(\check X)=\frac{(-1)^n}{\deg(\varphi)}
\int_X\frac{(1+H)^{r-1}}{C(N)}.
\tag{5.5.1}
\]
\end{SGAresult}
\begin{proof}
对（5.2.8）应用（5.3.3），其中取 $M=N$，即 $X\hookrightarrow\mathbf P^r$ 的\Term{normal-sheaf}{法层}。
\end{proof}
对于一个\Term{quasi-projective}{拟射影}\Term{smooth}{光滑}\Term{variety}{簇} $S$，将它的\Term{total-chern-class}{总 Chern 类}记为（参见 5.3）
\[
C(S)\overset{\mathrm{dfn}}=C(T_S),
\tag{5.5.2}
\]
其中 $T_S$ 表示 $S$ 的\Term{tangent-bundle}{切丛}。

\SourcePage{260}
\Term{exact-sequence}{正合列}（5.1.2）的\Term{dual}{对偶}表明
\[
C(\mathbf P^r)=(1+H)^{r+1}
\quad\text{在 }\operatorname{CH}(\mathbf P^r)\text{ 中}。
\tag{5.5.3}
\]
\Term{exact-sequence}{正合列}（5.1.1）的\Term{dual}{对偶}表明
\[
C(N)\,C(X)=C(\mathbf P^r)|_X
\quad\text{在 }\operatorname{CH}(X)\text{ 中}。
\tag{5.5.4}
\]
结合（5.5.3）和（5.5.4），（5.5.1）可改写为：
\begin{SGAresult}{推论}{5.6}
有公式
\[
\deg(\check X)=\frac{(-1)^r}{\deg(\varphi)}
\int_X\frac{C(X)}{(1+H)^2}.
\tag{5.6.1}
\]
\end{SGAresult}

\textbf{5.7。}对于\Term{projective}{射影}\Term{smooth}{光滑}\Term{variety}{簇} $S$，定义它的\Term{euler-poincare-characteristic}{Euler--Poincaré 示性数}为
\[
\chi(S)=\int_S C(S).
\tag{5.7.1}
\]
\begin{SGAresult}{命题}{5.7.2}
以 $Y$ 表示 $X$ 的一个\Term{smooth}{光滑}\Term{hyperplane-section}{超平面截面}，以 $\Delta$ 表示 $X$ 与 $\mathbf P^r$ 中\Term{codimension}{余维}为 $2$ 的一个\Term{linear-subspace}{线性子空间}相交所得的\Term{smooth}{光滑}\Term{section}{截面}。有
\[
\deg(\check X)=\frac{(-1)^n}{\deg(\varphi)}
\bigl[\chi(X)+\chi(\Delta)-2\chi(Y)\bigr].
\tag{5.7.3}
\]
\end{SGAresult}
\SourcePage{261}
\begin{proof}
以 $i:Y\hookrightarrow X$ 表示\Term{inclusion}{包含映射}。有
\[
N_{Y\hookrightarrow X}=H|_Y=i^*H,
\]
因此\Term{exact-sequence}{正合列}
\[
0\longrightarrow T_Y\longrightarrow i^*T_X
\longrightarrow i^*H\longrightarrow0
\]
给出
\[
C(Y)=i^*\left(\frac{C(X)}{1+H}\right).
\tag{5.7.4}
\]
所以
\[
\chi(Y)=\int_Yi^*\left(\frac{C(X)}{1+H}\right),
\]
应用\Term{projection-formula}{投影公式}便有
\[
\chi(Y)=\int_Xi_*i^*\left(\frac{C(X)}{1+H}\right)
=\int_X\frac{H\cdot C(X)}{1+H}.
\tag{5.7.5}
\]
重复这一计算，对 $\Delta$ 得到
\[
\begin{aligned}
\chi(\Delta)
&=\int_Y\frac{H\,C(Y)}{1+H}
=\int_Yi^*\left(\frac{H}{1+H}\cdot\frac{C(X)}{1+H}\right)\\
&=\int_X\frac{H^2C(X)}{(1+H)^2},
\end{aligned}
\tag{5.7.6}
\]
因而
\[
\begin{aligned}
\chi(X)-2\chi(Y)+\chi(\Delta)
&=\int_X C(X)\left[1-\frac{2H}{1+H}+\frac{H^2}{(1+H)^2}\right]\\
&=\int_X\frac{C(X)}{(1+H)^2}.
\end{aligned}
\]
结合 5.6.1，这就完成了证明。
\end{proof}

\SourcePage{262}
\SGAsection{6}{一般基底的情形}
\textbf{6.1。}设 $S$ 是一个\Term{scheme}{概形}，$\bar s$ 是 $S$ 的一个\Term{geometric-point}{几何点}，由某个 $s\in S$ 的\Term{residue-field}{剩余域} $k(s)$ 的一个\Term{algebraic-closure}{代数闭包}定义。以 $\mathbf P^r$ 表示 $S$ 上的 $r$ 维\Term{projective-space}{射影空间}，以 $\check{\mathbf P}^{r}$ 表示\Term{dual-projective-space}{对偶射影空间}，以 $I\subset\mathbf P^r\times_S\check{\mathbf P}^{r}$ 表示\Term{incidence-variety}{关联簇}。设 $D\subset\check{\mathbf P}^{r}$ 是一条直线，由适当的\Term{grassmannian}{Grassmann 簇}的一个 $S$ 值点定义；以 $\Delta$ 表示相应\Term{hyperplane}{超平面}\Term{pencil}{铅笔}的\Term{axis}{轴}（参见 2.1）。

设 $X$ 是一个\Term{proper}{固有}\Term{smooth}{光滑}的 $S$-\Term{scheme}{概形}，纯\Term{dimension}{维数}为 $n\geq1$，并赋有一个\Term{projective-embedding}{射影嵌入} $X\hookrightarrow\mathbf P^r$。令
$\widetilde X=(X\times_SD)\cap I$，这是 $D$ 上的\Term{scheme}{概形}，其\Term{fiber}{纤维}为 $D$ 所参数化的 $X$ 的\Term{hyperplane-section}{超平面截面}；以 $u$ 表示 $\widetilde X$ 到 $D$ 的\Term{projection}{投影}。

用下标 $\bar s$ 表示沿 $\bar s\longrightarrow S$ 的\Term{base-change}{基变换}。

假设 $\Delta_{\bar s}$ 与 $X_{\bar s}$ \Term{transverse}{横截}。则在 $s$ 在 $S$ 中的某个\Term{neighborhood}{邻域} $U$ 上，$\Delta$ 与 $X$ \Term{transverse}{横截}。以下假定 $U=S$。于是\Term{scheme}{概形} $\widetilde X$ 在 $S$ 上\Term{smooth}{光滑}，而 $\Delta$ 在 $\widetilde X$ 中的\Term{inverse-image}{逆像}的某个\Term{neighborhood}{邻域}在 $D$ 上\Term{smooth}{光滑}。

再假设属于\Term{pencil}{铅笔} $D_{\bar s}$ 的 $X_{\bar s}$ 的某个\Term{hyperplane-section}{超平面截面}是\Term{smooth}{光滑的}。则将 $S$ 换成 $s$ 的一个适当\Term{neighborhood}{邻域}之后，由\Term{jacobian-ideal}{Jacobi 理想} $J^{n-1}(\widetilde X/D)$ 定义的 $\widetilde X$ 的\Term{subscheme}{子概形} $F$ 在 $S$ 上是\Term{finite}{有限}\Term{flat}{平坦的}。

\SourcePage{263}
\textbf{6.2。}\Term{pencil}{铅笔} $D$ 的\Term{exceptional-locus}{例外集}是 $D/S$ 的\Term{relative-cartier-divisor}{相对 Cartier 除子} $E$，由在 $S$ 上\Term{finite}{有限}\Term{flat}{平坦}的 $\mathcal O_D$-\Term{module}{模} $u_*\mathcal O_F$ 定义。在 $D$ 上\Term{local-scheme}{局部}地，若
\[
\mathcal O_D^n\xrightarrow{v}\mathcal O_D^n
\longrightarrow u_*\mathcal O_F\longrightarrow0
\]
是 $u_*\mathcal O_F$ 的一个\Term{presentation}{表示}，则 $E$ 的方程为 $\det(v)=0$。$E$ 的形成与任意\Term{base-change}{基变换}相容。

$D_{\bar s}$ 的\Term{divisor}{除子} $E_{\bar s}$ 可写成 $\sum n(P)\cdot P$，其中 $n(P)$ 是像为 $P$ 的 $u$ 的各个\Term{critical-point}{临界点}处\Term{jacobian-ideal}{Jacobi 理想} $J^{n-1}(\widetilde X/D)$ 的\Term{colength}{余长度}之和。

\begin{SGAresult}{命题}{6.3}
若 $D_{\bar s}$ 是 $X_{\bar s}$ 的\Term{hyperplane-section}{超平面截面}的一个\Term{lefschetz-pencil}{Lefschetz 铅笔}，且 $\operatorname{car}k(s)\ne2$ 或 $n$ 为偶数，则 $s$ 在 $S$ 中存在一个\Term{neighborhood}{邻域} $U$，使得：

\textup{(a)} 对 $U$ 的每个\Term{geometric-point}{几何点} $\bar t$，$D_{\bar t}$ 都是 $X_{\bar t}$ 的\Term{hyperplane-section}{超平面截面}的一个\Term{lefschetz-pencil}{Lefschetz 铅笔}。

\textup{(b)} 在 $U$ 上，\Term{exceptional-locus}{例外集} $E$ 在 $S$ 上是\Term{etale}{平展的}。
\end{SGAresult}
\begin{proof}
为证明 (a)、(b)，只须注意，在所作假设下，$D_{\bar s}$ 满足 2.2 的条件 c)，当且仅当\Term{divisor}{除子} $E_{\bar s}$ 是\Term{reduced}{约化的}。援引 3.3，也可以只使用这一\Term{equivalence-relation}{等价关系}中较容易的一半（$\Longrightarrow$）。
\end{proof}

\textbf{6.3。}在上述例外情形中，这个命题不再成立。不过，若 $S$ 的\Term{characteristic}{特征}为 $2$ 且 $n$ 为奇数，则在 $S$ 中存在一个\Term{neighborhood}{邻域} $U$，使 (a) 在其上成立，并且可以定义一个在 $U$ 上\Term{etale}{平展}的\Term{divisor}{除子} $E'$，使得 $E=2E'$。我们不会使用这一结果。

\textbf{参考文献}
\begin{enumerate}
\item[{[1]}] A. Grothendieck，\emph{Chern 类理论}（La théorie des Classes de Chern），Bull. Soc. Math. France，86，1958，137--154。
\end{enumerate}
